\begin{proof}[Proof of \cref{obj:lem:ht-transfer}]
For the cited expected absolute-supremum contraction conclusion, we invoke \citep[Theorem 12(4), with the convention of Definition 2]{BartlettMendelson2002Complexities}.
The calculation below uses the measurable, centered, square-integrable prognosis supplied by \cref{obj:def:sobolev-class}. Since \(n,d\geq2\), all divisions by \(n\) and first-coordinate evaluations are well defined.

\begin{enumerate}
\item
Choose \(X\) in the probability-one set on which every fairness equality in \cref{obj:ass:fair} holds. This set is chosen independently of the schedule \(Y\). For either value of \(Z_i\),
\[
2Z_iY_i\!\left(\frac{Z_i+1}{2}\right)
=
Y_i(1)-Y_i(0)+2Z_iq_i.
\]
Summing gives the pointwise identity
\[
\widehat\theta
=
\theta_n+\frac2n\sum_{i=1}^n q_iZ_i.
\]
Under \cref{obj:ass:assignment-independence}, assignment conditional on \((X,Y)\) uses the kernel \(\pi(X)\). Consequently,
\[
\mathbb E_\pi[\widehat\theta\mid X,Y]
=
\theta_n+\frac2n\sum_{i=1}^n
q_i\mathbb E_\pi[Z_i\mid X]
=
\theta_n.
\]
Expanding the centered square and using fairness in the definition of the assignment covariance yields
\[
\begin{aligned}
\operatorname{Var}_\pi(\widehat\theta\mid X,Y)
&=
\frac4{n^2}
\mathbb E_\pi\!\left[
\left(\sum_{i=1}^n q_iZ_i\right)^2
\,\middle|\,X\right]\\
&=
\frac4{n^2}\sum_{i=1}^n\sum_{j=1}^n
q_i\mathbb E_\pi[Z_iZ_j\mid X]q_j\\
&=
\frac4{n^2}q^{\mathsf T}C_\pi(X)q.
\end{aligned}
\]
These calculations apply to every real schedule on the same probability-one set, proving both conditional assertions. % lean: CausalSmith.Experimentation.BivariateSobolevDesignCapacity.ht_conditional_transfer

\item
Define the Gaussian completion's treatment-effect function and its sample average by
\[
\begin{aligned}
\tau_{\mathrm G}(x)
&=
0.2+0.1\sqrt2\sin(2\pi x_1),
&&x\in[0,1]^d,\\
\overline\tau_{\mathrm G}(X)
&=
\frac1n\sum_{i=1}^n\tau_{\mathrm G}(X_i),
&&X\in([0,1]^d)^n.
\end{aligned}
\]
The completion in \cref{obj:synth_4} gives
\[
\theta_n
=
\frac1n\sum_{i=1}^n
\bigl(\tau_{\mathrm G}(X_i)+\varepsilon_{i1}-\varepsilon_{i0}\bigr),
\qquad
\mathbb E_Y[\theta_n\mid X]
=
\overline\tau_{\mathrm G}(X).
\]
Each first coordinate is uniform on \([0,1]\), and integration over a full sine period gives
\[
\int_0^1\sin(2\pi v_{\mathrm{int}})\,dv_{\mathrm{int}}
=
\left[-\frac{\cos(2\pi v_{\mathrm{int}})}{2\pi}\right]_0^1
=
0.
\]
Thus
\[
\mathbb E_X[\tau_{\mathrm G}(X_i)]=0.2,
\qquad
\theta
=
\mathbb E_{X,Y}[\theta_n]
=
\mathbb E_X[\overline\tau_{\mathrm G}(X)]
=
0.2.
\]
The frozen benchmark evaluates to
\[
V^*=(0.1)^2+4=4.01.
\]
% lean: CausalSmith.Experimentation.BivariateSobolevDesignCapacity.gaussianTheta_eq_tau0

\item
For every covariate array and sign vector, define
\[
S_m(X,Z)=\sum_{i=1}^n Z_i m(X_i),
\qquad
(X,Z)\in([0,1]^d)^n\times\{-1,1\}^n.
\]
Cauchy--Schwarz and \(Z_i^2=1\) give the pointwise bound
\[
S_m(X,Z)^2
\leq
n\sum_{i=1}^n m(X_i)^2.
\]
Because \(\pi(X)\) is a probability law,
\[
0\leq
\mathbb E_\pi[S_m(X,Z)^2\mid X]
\leq
n\sum_{i=1}^n m(X_i)^2.
\]
Integrating against the uniform sample law therefore gives
\[
\mathbb E_{X,Z}[S_m(X,Z)^2]
\leq
n^2\int_{[0,1]^d}m(x)^2\,dx,
\]
and hence
\[
0\leq\mathcal L_{n,d,s}(\pi,m)
\leq
4n\int_{[0,1]^d}m(x)^2\,dx
<\infty.
\]
The same bound ensures integrability of the conditional imbalance second moment. % lean: CausalSmith.Experimentation.BivariateSobolevDesignCapacity.loss_lt_top_of_memLp

\item
Write \(\varepsilon=(\varepsilon_{i0},\varepsilon_{i1})_{i=1}^n\) for the Gaussian noise array in \cref{obj:synth_4}, and define the selected-noise contribution on its full domain by
\[
N_{\mathrm{sel}}(Z,\varepsilon)
=
\frac2n\sum_{i=1}^n
Z_i\varepsilon_{i,(Z_i+1)/2},
\qquad
(Z,\varepsilon)\in\{-1,1\}^n\times\mathbb R^{n\times2}.
\]
Substituting the completion, for either assignment sign, gives
\[
2Z_iY_i\!\left(\frac{Z_i+1}{2}\right)
=
\tau_{\mathrm G}(X_i)
+2Z_i m(X_i)
+2Z_i\varepsilon_{i,(Z_i+1)/2}.
\]
Consequently,
\[
\widehat\theta
=
\overline\tau_{\mathrm G}(X)
+\frac2n S_m(X,Z)
+N_{\mathrm{sel}}(Z,\varepsilon).
\]
For fixed \(Z\), the signed selected errors are independent across units, each with mean zero and variance one. With \(\mathbb E_\varepsilon\) denoting expectation over the Gaussian noise at fixed \(X,Z\), it follows that
\[
\begin{aligned}
\mathbb E_\varepsilon[N_{\mathrm{sel}}(Z,\varepsilon)]
&=
\frac2n\sum_{i=1}^n
\mathbb E_\varepsilon[
Z_i\varepsilon_{i,(Z_i+1)/2}]
=0,\\
\operatorname{Var}_\varepsilon(N_{\mathrm{sel}}(Z,\varepsilon))
&=
\frac4{n^2}\sum_{i=1}^n1
=\frac4n,\\
\mathbb E_\varepsilon[N_{\mathrm{sel}}(Z,\varepsilon)^2]
&=\frac4n.
\end{aligned}
\]
Thus
\[
\mathbb E_\varepsilon[\widehat\theta]
=
\overline\tau_{\mathrm G}(X)+\frac2n S_m(X,Z),
\]
and expansion of the square gives
\[
\begin{aligned}
\mathbb E_\varepsilon[(\widehat\theta-0.2)^2]
&=
\left(
\overline\tau_{\mathrm G}(X)-0.2+\frac2n S_m(X,Z)
\right)^2\\
&\quad+
2\left(
\overline\tau_{\mathrm G}(X)-0.2+\frac2n S_m(X,Z)
\right)
\mathbb E_\varepsilon[N_{\mathrm{sel}}(Z,\varepsilon)]\\
&\quad+
\mathbb E_\varepsilon[N_{\mathrm{sel}}(Z,\varepsilon)^2]\\
&=
\left(
\overline\tau_{\mathrm G}(X)-0.2+\frac2n S_m(X,Z)
\right)^2+\frac4n.
\end{aligned}
\]
% lean: CausalSmith.Experimentation.BivariateSobolevDesignCapacity.ht_completion_noise_centered_second_moment

\item
For almost every \(X\), fairness gives
\[
\mathbb E_\pi[S_m(X,Z)\mid X]
=
\sum_{i=1}^n m(X_i)\mathbb E_\pi[Z_i\mid X]
=
0.
\]
The finite assignment space permits exchanging the assignment average and the Gaussian-noise integral. Averaging the preceding noise mean therefore yields
\[
\mathbb E_{Y,Z}[\widehat\theta\mid X]
=
\overline\tau_{\mathrm G}(X).
\]
For the centered second moment, expand
\[
\begin{aligned}
\left(
\overline\tau_{\mathrm G}(X)-0.2+\frac2n S_m(X,Z)
\right)^2
&=
(\overline\tau_{\mathrm G}(X)-0.2)^2\\
&\quad+
\frac4n(\overline\tau_{\mathrm G}(X)-0.2)S_m(X,Z)
+\frac4{n^2}S_m(X,Z)^2.
\end{aligned}
\]
The middle term has assignment mean zero, so
\[
\mathbb E_{Y,Z}[(\widehat\theta-0.2)^2\mid X]
=
(\overline\tau_{\mathrm G}(X)-0.2)^2
+\frac4{n^2}\mathbb E_\pi[S_m(X,Z)^2\mid X]
+\frac4n.
\]
The bounded treatment-effect function, the integrable imbalance bound, and Gaussian second moments ensure that this expression is integrable over \(X\), and that the actual estimator has a finite centered second moment. % lean: CausalSmith.Experimentation.BivariateSobolevDesignCapacity.ht_gaussian_kernel_centered_second_moment

\item
Integrating the conditional mean over the covariates gives
\[
\mathbb E_{X,Y,Z}[\widehat\theta]
=
\mathbb E_X[\overline\tau_{\mathrm G}(X)]
=
0.2.
\]
Square integrability and the iterated construction of the experiment law now imply
\[
\operatorname{Var}_{X,Y,Z}(\widehat\theta)
=
\mathbb E_{X,Y,Z}[(\widehat\theta-0.2)^2]
=
\mathbb E_X\!\left[
\mathbb E_{Y,Z}[(\widehat\theta-0.2)^2\mid X]
\right].
\]
% lean: CausalSmith.Experimentation.BivariateSobolevDesignCapacity.ht_gaussian_joint_mean

\item
The squared-sine integral is
\[
\begin{aligned}
\int_0^1\sin^2(2\pi v_{\mathrm{int}})\,dv_{\mathrm{int}}
&=
\int_0^1\frac{1-\cos(4\pi v_{\mathrm{int}})}2\,dv_{\mathrm{int}}\\
&=
\left[
\frac{v_{\mathrm{int}}}{2}
-\frac{\sin(4\pi v_{\mathrm{int}})}{8\pi}
\right]_0^1
=\frac12.
\end{aligned}
\]
Together with the treatment-effect mean already computed, this gives
\[
\operatorname{Var}_X(\tau_{\mathrm G}(X_i))
=
(0.1)^2\,2\int_0^1
\sin^2(2\pi v_{\mathrm{int}})\,dv_{\mathrm{int}}
=
(0.1)^2.
\]
Independence of the original covariate vectors consequently yields
\[
\begin{aligned}
\mathbb E_X[
(\overline\tau_{\mathrm G}(X)-0.2)^2]
&=
\operatorname{Var}_X(\overline\tau_{\mathrm G}(X))\\
&=
\frac1{n^2}\sum_{i=1}^n
\operatorname{Var}_X(\tau_{\mathrm G}(X_i))
=
\frac{(0.1)^2}{n}.
\end{aligned}
\]
By the definition of the finite signing loss,
\[
\frac4{n^2}
\mathbb E_X\!\left[
\mathbb E_\pi[S_m(X,Z)^2\mid X]
\right]
=
\frac{\mathcal L_{n,d,s}(\pi,m)}{n}.
\]
Integrating the conditional centered second moment thus gives
\[
\operatorname{Var}_{X,Y,Z}(\widehat\theta)
=
\frac{(0.1)^2}{n}
+\frac{\mathcal L_{n,d,s}(\pi,m)}{n}
+\frac4n.
\]
Multiplying by \(n\) and subtracting the evaluated benchmark proves
\[
n\operatorname{Var}_{X,Y,Z}(\widehat\theta)-V^*
=
\mathcal L_{n,d,s}(\pi,m).
\]
% lean: CausalSmith.Experimentation.BivariateSobolevDesignCapacity.ht_gaussian_outer_centered_second_moment
\end{enumerate}
\end{proof}
